% =============================================================================
% sample-accessible-template.tex
%
% A ready-to-use accessible LaTeX template.
% Every section below shows what this tool adds automatically — or what it
% checks and warns about — so you know what good accessible LaTeX looks like.
%
% QUICK START:
%   1. Fill in the CUSTOMIZE sections marked below.
%   2. Compile with: pdflatex sample-accessible-template.tex
%   3. Or let the tool handle a new file:
%        python3 latex-accessibility.py add yourfile.tex
% =============================================================================

\documentclass[12pt]{article}


% =============================================================================
% ACCESSIBILITY PACKAGES
% These are added automatically by: python3 latex-accessibility.py add yourfile.tex
%
% hyperxmp   — embeds XMP metadata stream in the PDF (PDF/UA clause 7.1)
% hyperref   — clickable links, PDF metadata, bookmarks (must come before bookmark)
% bookmark   — improved PDF navigation panel support
% enumitem   — accessible list formatting with consistent spacing
% =============================================================================
\usepackage{hyperxmp}
\usepackage{hyperref}
\usepackage{bookmark}
\usepackage{enumitem}


% =============================================================================
% PDF METADATA  [CUSTOMIZE: fill in your title, author, subject, and keywords]
%
% pdfdisplaydoctitle=true  — shows the pdftitle in the viewer title bar, not
%                            the filename (PDF/UA clause 7.1 — DisplayDocTitle)
% pdfuapart=1              — declares PDF/UA conformance in the XMP stream
%                            (PDF/UA clause 5 — conformance declaration)
% colorlinks/urlcolor      — coloured links instead of boxes; the RGB value
%                            below gives 7:1 contrast on white (WCAG AAA)
% =============================================================================
\hypersetup{
    pdfdisplaydoctitle=true,
    pdfuapart=1,
    colorlinks=true,
    urlcolor=[RGB]{0,40,150},
    linkcolor=[RGB]{0,40,150},
    citecolor=[RGB]{0,40,150},
    pdftitle={Your Document Title},          % CUSTOMIZE
    pdfauthor={Your Name},                   % CUSTOMIZE
    pdfsubject={Course or Subject},          % CUSTOMIZE
    pdflang={en-US},
    pdfkeywords={keyword1, keyword2}         % CUSTOMIZE
}


% =============================================================================
% PDF BOOKMARKS
% Added automatically by the tool.
%   numbered — bookmark entries include section numbers
%   open     — bookmark panel is open by default in PDF viewers
% =============================================================================
\bookmarksetup{
    numbered,
    open,
}


% =============================================================================
% LIST SPACING
% Added automatically by the tool.
% nosep removes extra vertical space inside lists for cleaner screen-reader output.
% =============================================================================
\setlist{nosep}


% =============================================================================
% TITLE BLOCK  [CUSTOMIZE]
% =============================================================================
\title{Your Document Title}
\author{
    Your Name \\
    Website: \href{https://example.com}{example.com} \\
    \date{\vspace{-5ex}}
}


% =============================================================================
% DOCUMENT BODY
% =============================================================================
\begin{document}

\maketitle


% =============================================================================
% ACCESSIBILITY NOTICE
% Added automatically by the tool when you provide an HTML URL.
% Replace the URL below with the actual address of the HTML version.
%
%   python3 latex-accessibility.py add yourfile.tex
%   (it will prompt for the HTML URL)
% =============================================================================
\section*{Accessibility Notice}
This document is also available in HTML format at:

Link: \href{https://example.com/yourdoc.html}{\textbf{https://example.com/yourdoc.html}}

The HTML version provides enhanced accessibility features including keyboard
navigation, screen reader support, responsive design, dark mode support, and
high contrast options.


% =============================================================================
% LINKS
% Wrap all URLs in \href{url}{\textbf{url}} with a "Link:" prefix.
% The full URL as link text is WCAG-compliant — it is self-describing.
% Plain \url{} is also acceptable but less visually distinct in print.
%
% The tool wraps bare URLs automatically.  If you write links by hand, use:
%   Link: \href{https://example.com}{\textbf{https://example.com}}
% =============================================================================
\section*{Introduction}

Example link to a resource:
Link: \href{https://example.com/resource}{\textbf{https://example.com/resource}}


% =============================================================================
% LISTS
% Use \begin{itemize} for unordered and \begin{enumerate} for ordered lists.
% The enumitem package (above) ensures consistent, accessible list structure.
% =============================================================================
\section*{Objectives}
\begin{itemize}
    \item First objective.
    \item Second objective.
\end{itemize}

\section*{Steps}
\begin{enumerate}
    \item First step.
    \item Second step, with sub-steps:
        \begin{enumerate}
            \item Sub-step A.
            \item Sub-step B.
        \end{enumerate}
    \item Third step.
\end{enumerate}


% =============================================================================
% FIGURES
% The tool checks for and warns about figures missing \caption.
% Always include:
%   1. An alt text comment — added automatically if missing:
%        % Alt text: [Add description here for accessibility]
%      Replace the placeholder with a real description before publishing.
%   2. A \caption{} — visible description; required for accessibility.
%
% Replace the filename and texts below with your own.
% =============================================================================
\section*{Example Figure}

\begin{figure}[h]
    \centering
    % Alt text: [REPLACE with a plain-English description of the image]
    % Replace the \fbox below with: \includegraphics[width=0.8\textwidth]{your-image-filename}
    \fbox{\parbox{0.75\textwidth}{\centering\vspace{1.5em}
        \texttt{[Replace this box with your \textbackslash includegraphics command]}
    \vspace{1.5em}}}
    \caption{A descriptive caption explaining what this figure shows.}
    \label{fig:example}
\end{figure}


% =============================================================================
% TABLES
% The tool checks for and warns about tables missing \caption.
% Always include:
%   1. A \caption{} above the tabular data.
%   2. A bold header row so column roles are clear.
% =============================================================================
\section*{Example Table}

\begin{table}[h]
    \centering
    \caption{An example table with an accessible caption.}
    \begin{tabular}{|l|l|l|}
        \hline
        \textbf{Column A} & \textbf{Column B} & \textbf{Column C} \\
        \hline
        Row 1, A & Row 1, B & Row 1, C \\
        Row 2, A & Row 2, B & Row 2, C \\
        \hline
    \end{tabular}
    \label{tab:example}
\end{table}


% =============================================================================
% COLOR
% Never use color as the ONLY way to convey information — colorblind readers
% will miss it.  The tool warns about \textcolor usage with no secondary cue.
%
% Instead of:  \textcolor{red}{Warning}
% Write:       \textcolor{red}{\textbf{Warning}}
%
% The bold/italic/underline cue carries the meaning for users who cannot see color.
% =============================================================================


% =============================================================================
% DELIVERABLES  [CUSTOMIZE]
% =============================================================================
\section*{Deliverables}

Describe what should be submitted here, including file names, formats, and
any screenshots or answers required.

\end{document}
